
\subsection{滤过 $A_\infty$ 双模的典范模型}
\label{subsec:5-4-2}
\liftIndex{I211}

现在讨论一个 $(C_1,C_0)$-滤过
$A_\infty$ 双模 $D$。
假设 $C_1,C_0,D$ 都有隙，
并且这里考虑的是以
$\Lambda_{0,\mathrm{nov}}$ 为系数的情形。
令
\[
\mathfrak n_{k_1,k_0}\colon
B_{k_1}C_1[1]\otimes D[1]\otimes
B_{k_0}C_0[1]\longrightarrow D[1]
\]
为 $D$ 的滤过 $A_\infty$ 双模结构。

由命题 5.2.43 可知，滤过
$A_\infty$ 双模的同伦范畴
只依赖于系数 $A_\infty$ 代数的同伦型；
由定理 5.4.2 可知，
每个 $A_\infty$ 代数都有典范模型。
因此，在以下讨论中，
可以假设 $C_1,C_0$ 是典范的。
下面的结果对应于定理 5.4.2。

\begin{theorem}{5.4.18}\label{thm:5-4-18}
存在 $(C_1,C_0)$ 上的有隙滤过
$A_\infty$ 双模 $(D',\mathfrak n')$，
它与 $D$ 同伦等价，且满足
\[
\mathfrak n'_{0,0}\equiv0
\pmod{\Lambda^+_{0,\mathrm{nov}}}.
\]
而且，同伦等价可取为有隙
$A_\infty$ 双模同态。
若 $C_1,C_0,D$ 都是同伦含幺的，
则 $D'$ 是含幺的。
\end{theorem}

证明与定理 5.4.2 的证明类似，
将留到第 5.4.4 小节给出。

\begin{definition}{5.4.19}\label{def:5-4-19}
若典范滤过 $A_\infty$ 代数上的
滤过 $A_\infty$ 双模 $D$
满足
$\mathfrak n_{0,0}\equiv0
\pmod{\Lambda^+_{0,\mathrm{nov}}}$，
就称 $D$ 为\emph{典范的}。
\end{definition}

\begin{proposition}{5.4.20}\label{prop:5-4-20}
两个典范滤过 $A_\infty$ 双模
之间的任意同伦等价都是同构。
\end{proposition}

证明与命题 5.4.5 相同，故从略。
将上述讨论应用于几何情形，
便得到下面的推论；
其证明将在第 5.4.6 小节给出。

\begin{corollary}{5.4.21}\label{cor:5-4-21}
对 $M$ 中干净相交的拉格朗日子流形
之任意相对 spin 对
$(L^{(0)},L^{(1)})$，
都可关联一个以
\[
\bigl(
C_{\mathrm{can}}(L^{(i)};\Lambda_{0,\mathrm{nov}}),
\mathfrak m_{\mathrm{can}}
\bigr),\qquad i=0,1,
\]
为系数代数的含幺典范滤过
$A_\infty$ 双模
\[
\bigl(
C_{\mathrm{can}}(L^{(1)},L^{(0)};
\Lambda_{0,\mathrm{nov}}),
\mathfrak n_{\mathrm{can}}
\bigr).
\]
在同构意义下，
该典范模型不依赖于相容殆复结构的选择，
也不依赖于定义过程中涉及的其他选择。
\end{corollary}

更具体地，
$C_{\mathrm{can}}(L^{(1)},L^{(0)};
\Lambda_{0,\mathrm{nov}})$
由下式给出。
\translatornote{原书实体页 345
在下式各直和项只写
$H^*(R_h;\Lambda_{0,\mathrm{nov}})$，
但同页紧接的次数说明
明确要求 $S$ 取值于
$\Theta^-_{R_h}$；
第 3.7.5 小节的博特--莫尔斯
双模构造也保留该定向局部系统。
此处补回它，见疑误 A-0317。}
\[
C_{\mathrm{can}}(L^{(1)},L^{(0)};
\Lambda_{0,\mathrm{nov}})
=
\bigoplus_{
[h,w]\in\widetilde\pi_0
(L^{(1)}\cap L^{(0)})}
H^*\bigl(
R_h;\Theta^-_{R_h}
\otimes_{\mathbb Q}\Lambda_{0,\mathrm{nov}}
\bigr).
\]
这里
\[
L^{(1)}\cap L^{(0)}
=
\bigcup_{h\in\pi_0(L^{(1)}\cap L^{(0)})}
R_h
\]
是按连通分支的分解。
关于次数，
若
\[
S\otimes[h,w]\otimes T^\lambda e^\mu
\in C_{\mathrm{can}}^k
(L^{(1)},L^{(0)};\Lambda_{0,\mathrm{nov}}),
\qquad
S\in C^\ell(R_h;\Theta^-_{R_h}),
\]
则
\[
k
=\deg\bigl(
S\otimes[h,w]\otimes T^\lambda e^\mu
\bigr)
=\ell+\mu([h,w])+2\mu.
\]
参见第 3.7.5 小节。

不过要注意，
$C_{\mathrm{can}}(L^{(1)},L^{(0)};
\Lambda_{0,\mathrm{nov}})$
并不是拉格朗日子流形的哈密顿同痕不变量。
如第 3.7.6 小节已经指出，
弗洛尔上同调在哈密顿同痕下的同构
只对 $\Lambda_{\mathrm{nov}}$ 系数成立，
对 $\Lambda_{0,\mathrm{nov}}$ 系数
并不成立。
第 6.5.3 小节将进一步讨论这一点。
对于以 $\Lambda_{\mathrm{nov}}$ 为系数的
滤过 $A_\infty$ 双模，
我们尚不知道合适的典范模型概念。

定理 F 立即由推论 5.4.21 推出。
